\paragraph{Problem 1 answer}\GradescopeComment{This should be on page 3 (top) of your PDF.}


\begin{enumerate}[label=(\alph*)]

\item[(a)] % Give a recurrence relation for $N(j)$ in terms of $N(1), \ldots, N(j-1)$.
  Here is the recurrence relation for $N(j)$: 

  \begin{blue}
    \[
      N(j) = \REPLACEME
    \]
  \end{blue}
  
  % For reference, here's the LaTeX source for the recurrence in Section 7.7 for counting pairwise-disjoint subsets.
  % \[
  %   N(j) = \begin{cases}
  %     1 & \text{ if } j=0 \\
  %     N(j-1) + N(i'-1), \text{ where }  i' = \min\{ i : e_i \ge s_j \}
  %     & \text{ if } j>0.
  %   \end{cases} 
  % \]

\item[(b)]  % Explain precisely in at most four lines why your recurrence relation is correct. 
  Here is why the above recurrence relation holds:

  \begin{blue}
    \REPLACEME
  \end{blue}
  
\item[(c)] % Give a recurrence relation for $M(j)$ in terms of $M(1), \ldots, M(j-1)$.
  Here is the recurrence relation for $M(j)$:

  \begin{blue}
    \[
      M(j) = \REPLACEME
    \]
  \end{blue}

  % For reference, here's the LaTeX source for the recurrence in Section 7.7 for weighted activity selection.
  % \[
  %   V(j) = \begin{cases}
  %     0 & \text{ if } j=0 \\
  %     \max\big( V(j-1), V(i'-1) + v_j\big), \text{ where }  i' = \min\{ i : e_i \ge s_j \},
  %     & \text{ if } j>0.
  %   \end{cases} 
  % \]

\item[(d)]  % Explain precisely in at most three or four lines why your recurrence relation is correct. 
  Here is why the above recurrence relation holds:

  \begin{blue}
    \REPLACEME
  \end{blue}
  
\end{enumerate}

\newpage

\begin{enumerate}

\item[(e)]   % What is the asymptotic worst-case running time of your algorithm in terms of $n$? 
  The asymptotic worst-case running time is \GradescopeComment{This should be on page 4 (top) of your PDF.}
  \BLUE{
    \[
      \Theta(\REPLACEME).
    \]
  }

\item[(f)] Here is why the above bound holds:
  % Explain precisely in at most three or four lines why the bound you give is correct. 

  \begin{blue}
    \REPLACEME
  \end{blue}
  
\end{enumerate}

\begin{itemize}

\item[(g)]  Here is the resulting iterative dynamic-programming algorithm:
  % Give the resulting iterative dynamic-programming algorithm, for computing $M(n)$.
  % Present your algorithm in pseudo-code,
  % following the guidelines in Lecture Note 1 and the examples in Lecture Note 7.
  %
  % For readability, use mathematical expressions (such as a sums, maximums, etc.)
  % that any competent programmer could implement in an obvious way,
  % without giving the details of how to compute those expressions.

  \begin{myframed}
    \begin{steps}
    \ITEM[]
      \textsf{freeLancer}$(\ell=(\ell_1, \ldots, \ell_n), h=(h_1, \ldots, h_n))$:
      \comment{--- iterative version}
      \begin{blue}
        
      \step \REPLACEME

      \step $\vdots$            % edit these steps as necessary

      \step return \REPLACEME

      \end{blue}
    \end{steps}
  \end{myframed}

  % For reference, here's the LaTeX source for the maxSchedule algorithm in Section 7.7:
  %
  % \begin{myframed}
  %   \begin{steps}
  %   \ITEM[]
  %     \textsf{maxSchedule}$(C=\{[s_1, e_1], [s_2, e_2], \ldots, [s_n, e_n]\}, v)$:
  %     \comment{--- iterative version}
  %     \step for $j=0,1,2,\ldots, n$:
  %     \step~~~set~\(W[j] = \begin{cases}
  %       0 & \text{ if } j=0 \\
  %       \max( W[j-1], W[i'-1] + v_j), \text{ where }  i' = \min\{ i : e_i \ge s_j \}
  %       & \text{ if } j>0.
  %     \end{cases}\)
  %     \step return $W[n]$
  %   \end{steps}
  % \end{myframed}


\end{itemize}

\vfill


\paragraph{Problem \arabic{problem} collaboration and citations}~\GradescopeComment{This should be on page 4 (bottom) of your PDF.}

I collaborated on this problem with the following people:
\begin{blue}
  \begin{itemize}
  \item \REPLACEME                    % write "none" if none
  \end{itemize}
\end{blue}

My answers use ideas from the following sources (other than the lecture notes and textbooks): 
\begin{blue}
  \begin{itemize}
  \item \REPLACEME                    % write "none" if none
  \end{itemize}
\end{blue}



%%% Local Variables:
%%% mode: latex
%%% TeX-master: "homework_8"
%%% End:
